\subsection{Partial Derivatives}

\content{
\begin{definition} If $f(x,y)$ is a function of two variables, then the partial
derivative of $f$ with respect to $x$ at a point $(a,b)$ is given by 
$$ f_x(a,b) = \lim_{h\rightarrow 0}\frac{f(a+h,b)-f(a,b)}{h}. $$
Similarly, the partial derivative of $f$ with respect to $y$ is given by 
$$ f_y(a,b) = \lim_{h\rightarrow 0}\frac{f(a,b+h)-f(a,b)}{h}. $$
When these derivatives exist. 
\end{definition}
}{% presentation
\vspace{1in}
}{% nopresentation
}{% comments
Note that there are several notations for the partial derivative of $f$:
$f_x(x,y)$, $\frac{\del f}{\del x}$, $D_x f$.
}

\content{
\begin{example} Find the first partial derivatives of the function 
$$ f(x,y) = x^3+x^2y^3-2y^2. $$
Find $f_x(2,1)$.
\end{example}
}{% presentation
\vspace{3in}
}{% nopresentation
}{% comments
}

\content{
\begin{example} Find the first partial derivatives of the function 
$$ f(x,y) = \sin\left(\frac{x}{1+y}\right). $$
\end{example}
}{% presentation
\vspace{3in}
}{% nopresentation
}{% comments
}

\content{
Interpretation of partial derivatives:
}{% presentation
\vspace{3in}
}{% nopresentation
}{% comments
Draw a surface and demonstrate what the partial derivatives mean in the context
of two variable functions. 
}

\subsubsection{Higher Derivatives}

\content{
\begin{example} Find the second partial derivatives of 
$$ f(x,y) = x^3+x^2y^3-2y^2. $$
\end{example}
}{% presentation
\vspace{4in}
}{% nopresentation
}{% comments
Note that there are several different notations for this as well:
$$ \frac{\del^2 f}{\del x \del y}\quad f_{xy}\quad D_{xy}. $$

Also note that $f_{xy}$ and $f_{yx}$ are equal in this example, the next page
introduces the theorem which states when this happens. 
}

\content{
\begin{theorem} Suppose that $f$ is defined on a disk $D$ that contains the
point $(a,b)$. If the functions $f_{xy}$ and $f_{yx}$ are both continous on $D$,
then 
$$ f_{xy}(a,b) = f_{yx}(a,b). $$
\end{theorem}
}{% presentation
}{% nopresentation
}{% comments
}

\subsection{Tangent Planes and Linear Approximations}
\content{

}{% presentation
\vspace{3in}
}{% nopresentation
}{% comments
Derive the equation of the tangent plane. This ends up being 
$$ z-z_0 = f_x(a,b)(x-a)+f_y(a,b)(y-b). $$
}

\content{
\begin{example} Find the tangent plane to the elliptic paraloloid 
$$ z = 2x^2+y^2. $$
\end{example} 
}{% presentation
\vspace{3in}
}{% nopresentation
}{% comments
}

\content{
The tangent plane can be used to approximate values of a function near a point
whose value is known:
\begin{example} Use the tangent plane to approximate $f(1.1,-0.1)$ where 
$$ f(x,y) = xe^{xy}. $$
\end{example}
}{% presentation
\vspace{3in}
}{% nopresentation
}{% comments
approx. $1.1e^{-0.11} \approx 0.98542$.
}

\subsubsection{Differentials}

\content{
Recall that for functions of a single variable, we have that the differential
$dy$ is given by 
$$ dy = f'(x)dx. $$
We have a similar definition for functions of two variables: 
\begin{definition} We define the differential $dz$, also called the total
differential, to be the value 
$$ dz = f_x(x,y)dx + f_y(x,y)dy=\frac{\del f}{\del x}dx + \frac{\del f}{\del y}dy. $$
\end{definition}
The Linearization of $f$ at $(a,b)$ can then be written 
$$ f(x,y) \approx f(a,b) + dz. $$
}{% presentation
}{% nopresentation
}{% comments
}

\content{
\begin{example} If 
$$ z = f(x,y) = x^2+3xy - y^2, $$
find the differential $dz$.
\end{example}
}{% presentation
\vspace{3in}
}{% nopresentation
}{% comments
}

\content{
\begin{example} The baes radius and height of a right circular cone are measured
as 10 cm and 25 cm, respectively, with a possible error in measurement of as
much as 0.01 cm in each. Use differentials to estimate the maximum error in the
calculated volume of the cone. 
\end{example}
}{% presentation
\vspace{4in}
}{% nopresentation
}{% comments
}
